\documentclass[10pt,a4paper]{amsart}
\usepackage[margin=19mm]{geometry}
\usepackage{amsmath,amssymb,mathtools}
\usepackage[scr=rsfs,cal=euler]{mathalfa}
\usepackage{xfrac}
\usepackage[unicode,hidelinks,bookmarks=false]{hyperref}
\newcommand{\ie}{i.e.\ }
\newcommand{\cf}{cf.\ }
\newcommand{\resp}{respectively\ }
\newcommand{\Subsection}{\subsection}
\providecommand{\nobreakdash}{\nobreak\hbox{-}}
\setlength{\emergencystretch}{2em}
\multlinegap0pt
\usepackage{amssymb}
\usepackage[shortlabels]{enumitem}
\usepackage{xcolor}
\newtheorem{lemma}{Lemma}
\usepackage{cleveref}
\crefname{lemma}{Lemma}{Lemmas}
\newcommand{\derivative}{\mathrm{d}}
\newcommand{\lie}[1]{\mathcal{L}_{#1}}
\DeclareMathOperator*{\ricci}{Ric}
\title{Ricci solitons from the perspectives of energy function}
\author{Shubham Yadav, Hemangi Madhusudan Shah}
\date{Source: arXiv:2606.14336v2 (CC BY 4.0)}
\begin{document}
\maketitle

\noindent\emph{Source and licence.} Ricci solitons from the perspectives of energy function. Version arXiv:2606.14336v2 (CC BY 4.0), distributed by arXiv under the Creative Commons Attribution 4.0 International licence, http://creativecommons.org/licenses/by/4.0/, as recorded in the arXiv licence metadata for this exact version and shown on the abstract page https://arxiv.org/abs/2606.14336v2. Full licence notice: \texttt{licenses/arxiv-2606.14336-CC-BY-4.0.txt}.\par\smallskip

\noindent\textit{Editorial scope.} Counted unit: the article's lemma lemma (source0.tex lines 1152--1159). The article's complete original proof of it is delivered with it (source0.tex lines 1160--1206, one \texttt{proof} environment of 47 lines). No mathematics is written, repaired or re-derived: the statement and the proof are the article's own lines, in the article's own notation, and no retained line is substituted or re-flowed.  No further source-local context is needed: every cross-reference of every retained line resolves inside the two delivered spans.  Nothing was omitted from any retained span, so the package carries no omission record.  No retained line contains a citation, so the wrapper carries no bibliography and no bibliography entry.  No retained line contains a figure environment, an image-inclusion command or drawing code, so no image is delivered and the package declares no asset dependency.  Editorial formatting only, in addition: an AMS \texttt{amsart} preamble that restates the article's own declarations and macros used by the retained lines, namely the environments \texttt{lemma} on separate counters, together with the standard packages those lines need.  The retained environments print in this document as Lemma 1 in order of appearance, whereas the article prints the same items with the numbering of its own full document.  The complete native source of the article as distributed in the arXiv e-print for this exact version is bundled unchanged as \texttt{sources/202777/source0.tex}.
\begin{lemma}
	Let $(M^n,g,V,\lambda)$ be a noncompact shrinking Ricci soliton with the bounded Ricci curvature. 
Then  there exists an $L_0$ such that for every 
	unit speed geodesic $\gamma: [0,L] \to M $ with $L\geq L_0$ and 
	$\gamma (0) = p \in M $,
	$$ E (L) \geq \frac{\lambda^2}{8} L^2 + \frac{\lambda}{2} a L + \frac{a^2}{2},  $$
	where $a$ is a constant that depends only on the geometry of the soliton on the unit ball $B(p,1)$.
\end{lemma}

\begin{proof}
	We first assume that $L\geq 2$.
	Let $E^{i}(p)$ denote an orthonormal basis along $\gamma$ with $E^n(t) = {\gamma}'(t)$ and 
   $E^{i}(t)$  denote  the parallel transport of $E^i$ over $\gamma(t)$.  Write
$V = V^i E_i$. Then for any Lipschitz function $\phi:[0,L] \to [0,1]$ with $\phi(0) = \phi(L)=0$, the second variation formula implies
	$$ (n-1) \int_0^L (\phi ' )^2 \, \derivative t 
		\geq \int_0^L \phi^2 \ricci (\gamma' , \gamma ' ) \, \derivative t. $$
	 Using the Ricci soliton equation and 
	$ \frac{1}{2} \lie{V} g (\gamma' , \gamma ' ) = \dot{V^n}, $
	in the above equation, we get
	$$ (n-1) \int_0^L (\phi ' )^2 \, \derivative t 
		\geq \frac{\lambda}{2} \int_0^L \phi^2 \, \derivative t 
		- \int_0^L \phi^2 \dot{V^n} \, \derivative t  . $$
	Subsequently,
	$$ (n-1) \int_0^L (\phi ' )^2 \, \derivative t 
		\geq \frac{\lambda}{2} \int_0^L \phi^2 \, \derivative t 
		- \int_0^L \dot{V^n} \, \derivative t		
		+ \int_0^L (1-\phi^2) \dot{V^n} \, \derivative t . $$
Using  the Ricci soliton equation again and noting the bound on Ricci curvature,
	\begin{align*}
		\int_0^L \left(1-\phi^2\right) \dot{V^n} \, \derivative t &= 
		\int_0^L \left(1-\phi^2\right) \left( \frac{\lambda}{2} - \ricci \left(\gamma' , \gamma'\right) \right) \, \derivative t \\
		&\geq \int_0^L \left(1-\phi^2\right) \left( \frac{\lambda}{2} - C \right) \, \derivative t. 
	\end{align*}
	Therefore,
	\begin{equation}\label{eqn: Lower estimate on Vn}
		(n-1) \int_0^L (\phi ' )^2 \, \derivative t 
		\geq - V^n(L) + V^n(0)		
		+ \int_0^L \left( \frac{\lambda}{2} - C \right)  \, \derivative t  
		+ C \int_0^L \phi^2 \, \derivative t .
	\end{equation}
	Now we choose,
	\begin{equation*}
		\phi = 
			\begin{cases}
				t, 	& 0 \leq t \leq 1 ; \\
				1, 	& 1 \leq t \leq L-1 ; \\
				L-t, 	& L-1 \leq t \leq L. \\
			\end{cases} 
	\end{equation*}
	Substituting in \cref{eqn: Lower estimate on Vn},
	$$ V^n (L) \geq V^n (0) - 2 (n-1) +\frac{\lambda}{2} L - \frac{4}{3} C , $$
	or
	$$ \sqrt{2E} (L) \geq V^n(L) \geq \frac{\lambda}{2} L + a, $$
	where $a$ is a constant that depends only on the geometry of the soliton on the unit ball $B(p,1)$. Since $\lambda>0$, we can choose $L_0$ large enough so that the right hand side of the above equation is positive. Then  we have for $L\geq L_0$, 
	$$ E (L) \geq \frac{\lambda^2}{8}L^2 + \frac{\lambda}{2} a L + \frac{a^2}{2}. $$
\end{proof}

\end{document}
